\documentclass[11pt, a4paper]{article}

% Packages
\usepackage[utf8]{inputenc}
\usepackage{amsmath, amssymb, amsfonts}
\usepackage{graphicx}
\usepackage{booktabs}
\usepackage{geometry}
\usepackage{subcaption}
\usepackage{float}
\usepackage{tabularx}
\usepackage{algorithm}
\usepackage{algpseudocode}
\usepackage{cite}
\usepackage{amsthm}
\theoremstyle{definition}
\usepackage{multirow}
\newtheorem{definition}{Definition}[section]

\geometry{top=2.5cm, bottom=2.5cm, left=2.5cm, right=2.5cm}

\title{\textbf{Causal Effect of Payment Protection Plans on Credit Default Risk}}
\author{Edoardo Lanzetti \and Valerio Viscovo}
\date{April 05, 2026}

\begin{document}

\maketitle

\begin{abstract}
Understanding whether Payment Protection Plans (PPPs) genuinely reduce credit default
risk is critical for financial institutions. However, evaluating observational data is
challenging because the treatment is not randomized. In this study, we estimate the
Average Causal Effect (ACE) of PPP enrollment on credit default. To rigorously control
for confounding, we compare a naive full-covariate baseline against a Post-Double-Selection
LASSO and two structural causal discovery methods (Rank PC and Fast-IAMB). The ACE is
estimated using Outcome Regression (G-computation) and Inverse Probability Weighting (IPW)
via logistic regression, with uncertainty quantified through a non-parametric bootstrap.
Our findings indicate a statistically significant positive causal effect of PPPs on default
risk, consistent with a mechanism of ``moral hazard'': protection from immediate financial
penalties may inadvertently encourage riskier repayment behavior. Furthermore, structurally
derived adjustment sets yield superior statistical efficiency compared to high-dimensional
predictive models, avoiding practical positivity violations that render high-dimensional IPW
estimates uninformative.
\end{abstract}

\section{Introduction and Causal Framework}

Banks routinely track whether customers will miss their credit card payments, as defaults
on repayments constitute the primary source of interest income for card issuers. To
mitigate the risk of default during financial hardship, certain customers may enroll in
Payment Protection Plans (PPPs), which allow them to pause repayments without incurring
immediate interest charges.

Understanding whether PPPs genuinely reduce default risk is of direct practical importance,
yet this question cannot be answered by a naive comparison of enrollees and non-enrollees.
Customers who elect to enroll in a PPP are likely to be at a higher \textit{a priori}
risk of missing payments, introducing a selection bias. Because of this structural
confounding, the true causal effect is distinct from the simple observational difference:
\begin{equation*}
    \text{Causal Effect} \neq \mathbb{E}[Y \mid A=1] - \mathbb{E}[Y \mid A=0]
\end{equation*}

To formally define the causal estimand, we rely on the Potential Outcomes framework
\cite{hernan2020causal}. Let the binary treatment $A \in \{0, 1\}$ denote PPP enrollment,
and the binary outcome $Y \in \{0, 1\}$ denote credit default. We define $Y^{a=1}$ and
$Y^{a=0}$ as the potential outcomes if a customer were compulsorily enrolled or not
enrolled, respectively. The Average Causal Effect (ACE) is defined as:
\begin{equation*}
    \text{ACE} = \mathbb{E}[Y^{a=1}] - \mathbb{E}[Y^{a=0}]
\end{equation*}

To identify the ACE from observational data, we rely on three core causal assumptions
\cite{hernan2020causal}:
\begin{enumerate}
    \item \textbf{Consistency:} The observed outcome matches the potential outcome for the
    treatment actually received: $Y = A Y^{a=1} + (1-A) Y^{a=0}$.
    \item \textbf{Conditional Exchangeability:} $Y^a \perp\!\!\!\perp A \mid \mathbf{L}$.
    Conditional on a set of observed covariates $\mathbf{L}$, the treatment assignment is
    independent of the potential outcomes. To satisfy this, $\mathbf{L}$ must block all
    \textit{backdoor paths} (spurious associations driven by common causes) without opening
    \textit{collider paths}. A collider $C$ is a common effect of two independent variables
    ($A \rightarrow C \leftarrow Y$); conditioning on $C$ induces a spurious dependence
    between $A$ and $Y$.
    \item \textbf{Positivity:} $0 < \mathbb{P}(A=1 \mid \mathbf{L} = \ell) < 1$ for all
    $\ell$. Every individual must have a strictly positive probability of receiving either
    treatment. Adjusting for high-dimensional sets $\mathbf{L}$ often violates this
    assumption in practice (practical positivity violation), leading to extreme estimation
    variance \cite{hernan2020causal}.
\end{enumerate}

This report addresses two questions: (1)~does PPP enrollment causally increase or decrease
the probability of credit default? (2)~how does the choice of adjustment set, informed by
causal discovery, affect both the magnitude and precision of the causal estimate? To answer
them, we compare four adjustment strategies of increasing theoretical sophistication---a
full covariate model, LASSO-based selection, and two causal discovery methods---and evaluate
each under two complementary estimators (OR and IPW) whose agreement serves as a robustness
check against model misspecification.

\section{Dataset Overview}

The dataset mimics a real credit repayment study ($n = 500$). The variables are summarized
in Table~\ref{tab:data-summary}. To ensure clarity in the subsequent causal graphs, we
specify the financial nomenclature:
\begin{itemize}
    \item \texttt{LIMIT\_BAL}: total credit granted (individual and family combined).
    \item \texttt{PAY\_0} to \texttt{PAY\_6}: repayment status for the previous six months
    on an ordinal scale ($-2 =$ no consumption, $-1 =$ paid in full, $0 =$ minimum payment,
    $1$ to $9 =$ months of delay).
    \item \texttt{BILL\_AMT1} to \texttt{BILL\_AMT6}: statement balance billed each month.
    \item \texttt{PAY\_AMT1} to \texttt{PAY\_AMT6}: actual amount paid each month.
\end{itemize}

\begin{table}[H]
    \centering
    \caption{Summary of dataset variables ($n = 500$).}
    \label{tab:data-summary}
    \begin{tabular}{p{5.5cm} p{3.5cm} p{5cm}}
        \toprule
        \textbf{Variable} & \textbf{Type} & \textbf{Summary / Range} \\
        \midrule
        Missed payment ($Y$)              & Binary           & 24\% defaulted \\
        PPP enrollment ($A$)              & Binary           & 44\% enrolled \\
        Credit limit (\texttt{LIMIT\_BAL})& Continuous       & $[10\,000;\ 650\,000]$ \\
        Age (\texttt{AGE})                & Continuous       & $[21;\ 73]$ \\
        Payment status (\texttt{PAY\_0--6})& Ordered categ.  & $-2$ to $9$ (integer) \\
        Bill amount (\texttt{BILL\_AMT1--6})& Continuous     & $[-2{,}650;\ 630{,}458]$ \\
        Amount paid (\texttt{PAY\_AMT1--6}) & Continuous     & $[0;\ 193{,}358]$ \\
        Education                         & Categorical      & Grad., Univ., High School \\
        \bottomrule
    \end{tabular}
\end{table}

\paragraph{Preprocessing.}
The ordinal \texttt{PAY} variables are encoded as integers preserving their natural order
(higher values indicate longer payment delay). The categorical \texttt{EDUCATION} variable
is dummy-encoded into three binary indicators (Graduate school as the reference category)
for use in the regression models. Two versions of the dataset are maintained throughout:
one with the original ordinal encoding for causal discovery algorithms (which require
numeric or factor inputs), and one with dummy-encoded education for the logistic regression
models used in estimation.

\section{Causal Discovery and Adjustment Strategies}

Identifying the correct adjustment set $\mathbf{L}$ is the central methodological challenge:
conditioning on too few variables leaves residual confounding, while conditioning on too
many risks positivity violations and efficiency loss. We evaluate four strategies of
increasing theoretical sophistication. Structural methods (Rank PC, Fast-IAMB) recover the
Markov blanket of $Y$ from the data-generating mechanism directly, providing a principled
answer to which variables must be adjusted for. The data-adaptive LASSO provides a flexible
benchmark that does not require distributional assumptions but lacks the graph-theoretic
guarantees of structural methods. The Full model serves as a naive upper bound. All analyses
use a fixed random seed for reproducibility.

\begin{definition}[Markov blanket]
Let $V$ be the set of all variables in a domain. A Markov Blanket of a target variable
$T \in V$, denoted $B(T) \subseteq V \setminus \{T\}$, is any subset of variables such
that $T$ is conditionally independent of all remaining variables given $B(T)$
\cite{pearl2009causality}:
\begin{equation*}
    T \perp\!\!\!\perp (V \setminus (B(T) \cup \{T\})) \mid B(T)
\end{equation*}
\end{definition}
A minimal Markov Blanket is called a \textit{Markov Boundary}, corresponding uniquely to
the parents, children, and spouses of $T$ in a DAG. Adjusting for the Markov Boundary of
$Y$ (excluding the treatment $A$) is sufficient to block all backdoor paths without
over-adjusting.

\subsection{Causal Structure via Rank PC Algorithm}

Standard implementations of the PC algorithm test conditional independence via partial
correlations, which assume the data follows a multivariate normal (Gaussian) distribution.
Financial data heavily violates this assumption due to extreme outliers (heavy tails) and
discrete categories. To overcome this, we utilize the Rank PC algorithm \cite{harris2013pc}.

Instead of analyzing raw data values, the Rank PC algorithm computes an empirical
Spearman rank correlation matrix ($\hat{\mathbf{R}}$). Unlike standard Pearson
correlation, which measures the linear relationship between exact numerical values,
Spearman's correlation evaluates the monotonic relationship between the relative rankings
of the data points (e.g., 1st, 2nd, 3rd largest). This rank-based approach makes it
highly robust to extreme financial outliers. We then apply a mathematical transformation
to this matrix:
\begin{equation*}
    \hat{\Sigma}_{ij} = 2 \sin\!\left(\frac{\pi}{6}\, \hat{R}_{ij}\right)
\end{equation*}
This transformation acts as a Gaussian copula, effectively mapping the rank correlations
into a latent Gaussian space. This allows the standard PC conditional independence tests
to remain valid without requiring the underlying data to be normally distributed.

The algorithm then iteratively removes and orients edges (using a significance level of
$\alpha = 0.05$), yielding a Completed Partially Directed Acyclic Graph (CPDAG). Finally,
our adjustment set is defined as the Markov Boundary of the outcome $Y$ on this graph
(its parents, children, and spouses), explicitly excluding the treatment variable $A$.

\subsection{Markov Blanket Discovery via Fast-IAMB}

Rather than learning the full graph, Fast-IAMB \cite{yaramakala2005speculative} directly
targets the Markov Boundary of $Y$, which is sufficient for adjustment without the
computational cost of global structure learning. To handle the mixed continuous-categorical
nature of the data and avoid distributional assumptions, all continuous variables are
discretized into three equiprobable categories (low, medium, high) using the quantile
method prior to running the algorithm.

Conditional dependence is evaluated via the $G^2$ likelihood-ratio test (equivalent to the
mutual information test for discrete data):
\begin{equation*}
    G^2 = 2 \sum_{i, j, k} O_{ijk} \ln \left( \frac{O_{ijk}}{E_{ijk}} \right)
\end{equation*}
where $i$, $j$, $k$ index states of the candidate variable, the target $T$, and the
current blanket $B(T)$ respectively; $O_{ijk}$ is the observed count and $E_{ijk}$ the
expected count under conditional independence. The algorithm proceeds in two phases:

\begin{algorithm}[H]
\caption{Fast-IAMB \cite{yaramakala2005speculative}}
\begin{algorithmic}[1]
\State \textbf{Input:} Target $T$, Variables $V$, Dataset $D$
\State \textbf{Phase I: Forward (Speculative) Phase}
\State Initialize $B(T) \gets \emptyset$
\Repeat
    \State Find $X \in V \setminus B(T)$ maximizing the $G^2$ statistic with $T$ given $B(T)$
    \If{$X \not\perp\!\!\!\perp T \mid B(T)$ (via $G^2$ test)}
        \State $B(T) \gets B(T) \cup \{X\}$
    \EndIf
\Until{No variable is added}
\State \textbf{Phase II: Backward (Elimination) Phase}
\For{each $X \in B(T)$}
    \If{$X \perp\!\!\!\perp T \mid B(T) \setminus \{X\}$ (via $G^2$ test)}
        \State $B(T) \gets B(T) \setminus \{X\}$
    \EndIf
\EndFor
\State \textbf{Return:} $B(T)$
\end{algorithmic}
\end{algorithm}

\subsection{Data-Adaptive Selection via LASSO}

Standard LASSO variable selection applied only to the outcome model can introduce omitted
variable bias: if a variable is a weak predictor of $Y$ but a strong driver of $A$, it
constitutes a confounder that LASSO may drop, leaving an open backdoor path. The
Post-Double-Selection procedure \cite{belloni2014inference} prevents this by fitting two
separate $\ell_1$-penalized logistic regression models---one for $Y$ and one for $A$---and
taking the union of the selected variables. For a generic binary target $T \in \{Y, A\}$:
\begin{equation*}
    \hat{\boldsymbol{\beta}}^{(T)} = \arg\min_{\boldsymbol{\beta}}
    \left( -\frac{1}{n} \sum_{i=1}^{n} \log P(T_i \mid \mathbf{X}_i, \boldsymbol{\beta})
    + \lambda_T \|\boldsymbol{\beta}\|_1 \right)
\end{equation*}
where $\lambda_T$ is selected by 10-fold cross-validation to minimise held-out deviance. The adjustment set is
$\mathbf{L} = \hat{S}_Y \cup \hat{S}_A$, where $\hat{S}_T = \{j : \hat{\beta}_j^{(T)} \neq 0\}$.

\section{Estimation and Inference}

Given an adjustment set $\mathbf{L}$, the ACE is estimated via logistic regression. For a
binary target $Z \in \{0,1\}$ and covariates $\mathbf{X}$, we model
$\text{logit}(\mathbb{P}(Z=1 \mid \mathbf{X})) = \boldsymbol{\beta}^\top \mathbf{X}$ and
obtain $\hat{\boldsymbol{\beta}}$ by maximum likelihood. We use two complementary
estimators; agreement between them strengthens confidence that the result is not an
artefact of one particular model:

\paragraph{1. Outcome Regression (G-computation):}
Let $Q(\mathbf{L}_i, a; \hat{\boldsymbol{\beta}}) = \mathbb{P}(Y=1 \mid A=a, \mathbf{L}_i;
\hat{\boldsymbol{\beta}})$ be the predicted default probability. G-computation marginalises
over the empirical covariate distribution:
\begin{equation*}
    \widehat{\text{ACE}}_{\text{OR}} = \frac{1}{n} \sum_{i=1}^n
    \left( Q(\mathbf{L}_i, 1; \hat{\boldsymbol{\beta}}) -
           Q(\mathbf{L}_i, 0; \hat{\boldsymbol{\beta}}) \right)
\end{equation*}
This estimator is consistent if the outcome model is correctly specified.

\paragraph{2. Inverse Probability Weighting (IPW):}
Let $\hat{\pi}(\mathbf{L}_i) = \mathbb{P}(A=1 \mid \mathbf{L}_i; \hat{\boldsymbol{\beta}})$
be the estimated propensity score, truncated to $[0.001, 0.999]$ to prevent numerical
instability from near-zero denominators:
\begin{equation*}
    \widehat{\text{ACE}}_{\text{IPW}} = \frac{1}{n} \sum_{i=1}^n
    \left( \frac{A_i Y_i}{\hat{\pi}(\mathbf{L}_i)} -
           \frac{(1-A_i)Y_i}{1-\hat{\pi}(\mathbf{L}_i)} \right)
\end{equation*}
This estimator is consistent if the propensity score model is correctly specified. When
$\mathbf{L}$ is high-dimensional, extreme propensity scores inflate the variance of the
IPW estimator, directly reflecting practical positivity violations.

\subsection{Uncertainty Quantification}

Statistical uncertainty is quantified via a non-parametric bootstrap ($B = 500$ resamples).
In each iteration, the outcome and propensity score models are refitted on the resampled
data for all four adjustment sets. The adjustment sets themselves are held fixed across
iterations. For the LASSO strategy, this is theoretically justified by the post-selection
inference guarantees of \cite{belloni2014inference}, which ensure that valid inference on
the treatment effect can be conducted conditional on the selected set. For the structural
methods (Rank PC and Fast-IAMB), the graph is fixed because structure learning algorithms
are not designed to be bootstrapped: the Rank PC algorithm is consistent for a fixed
population distribution \cite{harris2013pc}, but re-running it on each resample would
conflate the variance of the causal estimators with the instability of graph estimation
under small samples, making the resulting CIs difficult to interpret. Holding the graph
fixed therefore isolates the quantity of interest, the uncertainty in the ACE given a
fixed causal structure, at the cost of not propagating graph uncertainty into the
intervals. Note that for the LASSO strategy, holding the set fixed may also underestimate
total uncertainty given the moderate sample size ($n = 500$), as variable selection itself
is a source of randomness. The 95\% confidence intervals are derived via the percentile
method.

\section{Results and Discussion}

A naive comparison of default rates yields
$\hat{\Delta}_{\text{naive}} = \mathbb{E}[Y \mid A=1] - \mathbb{E}[Y \mid A=0] \approx 0.132$,
confirming substantial selection bias: customers who enroll in a PPP are observationally
more likely to default, but this raw difference conflates self-selection with any genuine
causal effect and cannot answer our first research question.

The adjustment sets identified by each strategy are detailed in Table~\ref{tab:adj-sets},
and the CPDAG recovered by Rank PC is shown in Figure~\ref{fig:cpdag}. The payment status
variables form a densely connected temporal subgraph, but only \texttt{PAY\_0} has a direct
edge into $Y$; the others influence $Y$ solely through \texttt{PAY\_0}. This
Markov-shielding property is why the Rank PC adjustment set retains only \texttt{PAY\_0}
from the payment history block, together with \texttt{LIMIT\_BAL} and \texttt{PAY\_AMT2}
(children of $Y$ in the CPDAG)and \texttt{BILL\_AMT3} and \texttt{PAY\_AMT6} (spouses of $Y$)---five variables in total. Fast-IAMB recovers an even sparser boundary of two variables (\texttt{PAY\_0} and
\texttt{PAY\_AMT4}), while LASSO selects 16 variables by design agnostic to the causal graph.

\begin{table}[ht]
    \centering
    \caption{Adjustment sets selected by each strategy. The Full Model (not shown) uses
    all 24 covariates. Structural methods identify the Markov Boundary of $Y$;
    LASSO takes the union of variables selected for $Y$ and $A$ separately.}
    \label{tab:adj-sets}
    \small
    \begin{tabular}{l l l}
    \toprule
    \textbf{LASSO (15 variables)} & \textbf{Rank PC (5 variables)} &
    \textbf{Fast-IAMB (2 variables)} \\
    \midrule
    \texttt{LIMIT\_BAL, AGE, SEX}            & \texttt{LIMIT\_BAL} & \texttt{PAY\_0}    \\
    \texttt{PAY\_0, PAY\_2, PAY\_3, PAY\_6}  & \texttt{PAY\_0}     & \texttt{PAY\_AMT4} \\
    \texttt{BILL\_AMT6}                      & \texttt{PAY\_AMT2}  &                    \\
    \texttt{PAY\_AMT1} to \texttt{PAY\_AMT6} & \texttt{BILL\_AMT3} &                    \\
    \texttt{EDU\_University}                 & \texttt{PAY\_AMT6}  &                    \\
    \bottomrule
\end{tabular}
\end{table}

\begin{figure}[H]
    \centering
    \includegraphics[width=1\textwidth]{cpdag.png}
    \caption{CPDAG recovered by the Rank PC algorithm ($\alpha = 0.05$, nonparanormal
    transformation). Solid arrows are directed edges; dashed lines are undirected edges
    whose orientation could not be determined from the data. Nodes are colour-coded by
    variable type. The Markov blanket of $Y$ 
    (parents \texttt{PAY\_0}, \texttt{PPP}; children \texttt{LIMIT\_BAL}, \texttt{PAY\_AMT2}; spouses \texttt{BILL\_AMT3}, \texttt{PAY\_AMT6}) is
    visually highlighted.}
    \label{fig:cpdag}
\end{figure}

Point estimates and bootstrap 95\% CIs for all eight estimator--adjustment combinations are
in Table~\ref{tab:bootstrap-results} and Figure~\ref{fig:bootstrap}.

\begin{figure}[H]
    \centering
    \includegraphics[width=0.82\textwidth]{ace_bootstrap_CI_.png}
    \caption{Point estimates (dots) and 95\% bootstrap percentile confidence intervals
    (bars) for the ACE of PPP enrollment on default probability, across four adjustment
    sets and two estimators. The dashed horizontal line marks ACE $= 0$. OR = Outcome
    Regression (G-computation); IPW = Inverse Probability Weighting. $B = 500$ resamples.}
    \label{fig:bootstrap}
\end{figure}

\begin{table}[H]
    \centering
    \caption{Bootstrap results ($B = 500$ resamples, percentile method). Width = upper
    bound $-$ lower bound of the 95\% CI. Intervals crossing zero are shown in bold.}
    \label{tab:bootstrap-results}
    \small
    \begin{tabular}{llccc}
        \toprule
        \textbf{Estimator} & \textbf{Adjustment} & \textbf{ACE} &
        \textbf{95\,\% CI} & \textbf{Width} \\
        \midrule
        \multirow{4}{*}{OR}
                             & Full      & 0.098 & $[0.017,\ 0.185]$          & 0.168 \\
                             & LASSO     & 0.105 & $[0.020,\ 0.189]$          & 0.169 \\
                             & Rank PC   & 0.118 & $[0.051,\ 0.187]$          & 0.136 \\
                             & Fast-IAMB & 0.111 & $[0.045,\ 0.181]$          & 0.136 \\
        \midrule
        \multirow{4}{*}{IPW}
                             & Full      & 0.087 & $\mathbf{[-0.051,\ 0.216]}$& 0.267 \\
                             & LASSO     & 0.084 & $\mathbf{[-0.028,\ 0.193]}$& 0.221 \\
                             & Rank PC   & 0.130 & $[0.062,\ 0.200]$          & 0.138 \\
                             & Fast-IAMB & 0.111 & $[0.041,\ 0.183]$          & 0.142 \\
        \bottomrule
    \end{tabular}
\end{table}

\subsection{Positivity and Variance Inflation}

The results expose a fundamental tension between predictive accuracy and causal validity.
Under IPW, the Full and LASSO adjustment sets yield CIs that cross zero
(Table~\ref{tab:bootstrap-results}), rendering the effect statistically indistinguishable
from null despite positive point estimates. The mechanism is a practical positivity
violation: conditioning on 24 (Full) or 16 (LASSO) variables allows the propensity model
to predict treatment assignment with high accuracy for some individuals, driving propensity
scores close to 0 or 1 and assigning them extreme inverse weights. This is not a failure of
IPW as a method, but a signal that the adjustment set is too rich relative to $n = 500$.
The Rank PC and Fast-IAMB sets (5 and 2 variables respectively) are sufficient to block the
relevant backdoor paths while preserving well-behaved propensity scores, yielding CI widths
of 0.136--0.142 under both estimators, roughly half those of the Full model. This directly
answers our second research question: causal discovery does not merely provide a theoretically
principled adjustment set, it also yields strictly more efficient estimates by preserving the
positivity assumption.

\subsection{Causal Interpretation}

Across all correctly specified structural models, the ACE is consistently positive and
statistically significant. Adjusting for the estimated confounding bias of approximately
3.4\,p.p.\ ($\hat{\Delta}_{\text{naive}} - \widehat{\text{ACE}}_{\text{OR, Full}} \approx
0.034$), the pure causal effect of PPP enrollment on default probability ranges from
approximately $+10$ to $+13$ percentage points, directly answering our first research
question: PPPs increase, rather than decrease, the probability of default.

These findings are consistent with a moral hazard mechanism \cite{arrow1963uncertainty}:
when customers are shielded from the immediate financial cost of missing a payment, the
incentive to manage liquidity carefully is weakened, ultimately increasing default
probability. Stiglitz and Weiss \cite{stiglitz1981credit} formalised this logic in the
credit market, showing that reducing the effective cost of non-repayment shifts borrowers
toward riskier strategies. The present results are consistent with this interpretation,
though the simulated nature of the dataset precludes definitive causal attribution to
this specific behavioral channel.

\section{Conclusion}

This study provides causal evidence that Payment Protection Plans increase the probability
of credit default. After controlling for confounding through structural causal discovery,
the estimated ACE ranges from $+11$ to $+13$ p.p., with CIs excluding zero across all
correctly specified adjustment sets.

Furthermore, causal discovery methods (Rank PC, Fast-IAMB)
recover parsimonious adjustment sets that satisfy the positivity assumption and yield
confidence intervals roughly half as wide as those from high-dimensional predictive
adjustment. In observational causal inference, parsimony in covariate selection is a requirement for valid inference.

\subsection{Limitations}

Four limitations deserve mention. First, the dataset contains only $n = 500$ observations,
limiting statistical power and making high-dimensional adjustment fragile. Second, causal
identification rests on conditional exchangeability (no unmeasured confounders); variables
such as individual risk tolerance are unobserved and, if correlated with both PPP enrollment
and default propensity, would bias the ACE upward---reinforcing but potentially
overstating the positive effect found here. Third, the Rank PC graph is sensitive to
$\alpha$ and the specific sample; graph uncertainty is not propagated into the bootstrap
CIs. Fourth, both estimators assume a logistic outcome model (linearity in the log-odds),
which was not formally validated; misspecification could bias the OR estimates while leaving
IPW estimates consistent, so the convergence between OR and IPW across structural adjustment
sets provides partial reassurance.

\subsection{Practical Implications}

These findings suggest that PPP enrollment should be treated as a risk indicator rather than
a mitigant. Financial institutions may benefit from segmenting customers before enrollment,
identifying profiles for which the behavioral response to payment deferral is more
pronounced, and pricing protection plans to reflect the increased default risk they may
induce.

\bibliographystyle{plain}
\bibliography{references}

\end{document}
